% !TEX root = ../../../main.tex
% Appendix entry: preprocessing.beads -- "BEADS baseline" (Appendix A).
\entryfield{Algorithm}
\code{_run} in \code{Adapters/preprocessing_beads.py} calls \code{remove_baseline} in
\code{Working/Preprocessing/beads.py}, passes on its \code{corrected} trace and its \code{meta},
and adds a list of plain-language \code{rules} to the meta. The core has two layers: a faithful
algorithm (\code{beads}) and a wrapper that adapts it to long, slow recordings
(\code{remove_baseline}).

\paragraph{The algorithm (\texttt{beads}).} BEADS \citep{ning2014chromatogram} models a trace as
$y = f + x + w$: $f$ a low-pass baseline, $x$ a sparse layer of events whose first and second
differences are also sparse, $w$ noise. It minimises
\[
F(x) = \tfrac12\lVert H(y-x)\rVert^2 + \lambda_0\sum_n \theta_\varepsilon(x_n; r)
  + \lambda_1\sum_n \phi([D_1x]_n) + \lambda_2\sum_n \phi([D_2x]_n),
\]
where $H = BA^{-1}$ is a zero-phase high-pass filter built from banded Toeplitz matrices of order
$d$ and cut-off $f_c$ (cycles per sample; \code{_filter_matrices}, the \texttt{BAfilt} of the
reference code), $\theta_\varepsilon$ is an asymmetric smoothed absolute value (positive $x$
costs 1 per unit, negative $x$ costs $r$), $\phi(v) = \sqrt{v^2+\varepsilon_1}$, and $D_1$,
$D_2$ are first and second differences. It is solved by majorisation--minimisation: each
iteration forms the diagonal weights $\Gamma = \mathrm{diag}\bigl((1+r)/(4\max(|x_n|,
\varepsilon_0))\bigr)$ and $\Lambda = \lambda_1 D_1^{\mathsf T}W_1D_1 + \lambda_2
D_2^{\mathsf T}W_2D_2$ with $W_k = \mathrm{diag}(1/\phi([D_kx]_n))$, then solves the banded
symmetric system
\[
x \leftarrow A\,\bigl(B^{\mathsf T}B + A^{\mathsf T}(\Lambda + 2\lambda_0\Gamma)A\bigr)^{-1}
  \bigl(B^{\mathsf T}BA^{-1}y - \lambda_0 A b\bigr), \qquad b_n = (1-r)/2,
\]
by banded Cholesky (\code{scipy.linalg.solveh_banded}), falling back to banded LU. It stops when
the cost changes by at most a fraction \code{tol} $= 10^{-3}$ between iterations, or after
\code{max_iter} $= 30$ iterations. The baseline is $f = y - x - H(y-x)$. The update equations
follow Laurent Duval's BSD-3 MATLAB \texttt{beads.m} line for line (with its \texttt{L1\_v2}
penalty, pybaselines' default); the paper's own text could not be read, so the equations are
the reference implementation's, and \unverified{their match to the published equations} is not
checked.

\paragraph{The wrapper (\texttt{remove\_baseline}).} Every step here is the project's adaptation,
not the paper's:
\begin{enumerate}
\item Check the inputs: at least 64 samples (\code{MIN_SAMPLES}) and 64 finite ones;
  $P = \mathit{baseline\_period\_s}\cdot f_s > 2$ samples; then $f_c = 1/P$ cycles per sample.
\item \emph{Missing samples.} Trim NaN runs at either end from the solve; bridge interior NaNs
  with straight lines.
\item \emph{Noise level.} $\sigma = \mathrm{MAD}(\Delta y)/(0.6745\sqrt2)$ on the finite samples
  (\code{robust_sigma}), falling back to the standard deviation of the difference, then to 0 for
  a straight line.
\item \emph{Polarity and weights.} \emph{drop} negates the trace and uses $r = 6$, \emph{spike}
  uses $r = 6$, \emph{both} uses $r = 1$ (symmetric). $(\lambda_0, \lambda_1, \lambda_2) =
  \mathit{sparsity}\times(0.15, 1.5, 1.2)$: the worked ratio $0.5 : 5 : 4$ times 0.3
  (\code{LAMBDA_RATIOS}); $\lambda_1$ and $\lambda_2$ are not exposed because small changes in
  $\lambda_1$ are reported to destabilise BEADS.
\item \emph{A coarse grid.} Average blocks of $q = \max\bigl(1, \min(\lfloor 0.02/f_c\rfloor,
  \lfloor m/64\rfloor)\bigr)$ samples (\code{grid_factor}, $m$ the trimmed length), lifting the
  cut-off to about 0.02 cycles per grid point (\code{FC_GRID}) where the banded system is well
  conditioned. At the defaults on a 1\,Hz recording $q = 144$ (144\,s blocks).
\item \emph{Ends.} Pad the grid at each end by one baseline period (at most the grid's length
  minus one), point-reflected about the end value ($2g_0 - g_k$), so level and slope continue;
  subtract a line through the medians of the two end regions (each the shorter of 5\% of the
  padded grid and a quarter period).
\item \emph{Scale.} Divide by $\sigma_g = \max\bigl(\mathrm{MAD}(Hz)/0.6745,\ \sigma/\sqrt q\bigr)$,
  the robust level of the high-passed grid measured without the padding (\code{grid_scale}),
  floored at the white-noise level. This makes \emph{sparsity} unit-free.
\item Run \code{beads} on the scaled grid with $f_c q$, order $d$ = \emph{filter\_order},
  $\varepsilon_0 = \varepsilon_1 = 10^{-4}$ (\code{EPS}).
\item Undo the scale, sign and end line; drop the padding; interpolate the baseline linearly
  back to every sample (continued linearly, not held flat, over the half block at each end);
  put every NaN back; return $\mathit{corrected} = y - \mathit{baseline}$.
\item Warn (in \code{meta["warnings"]}) when the event layer exceeds 3 scale units on more than
  80\% of the grid (not for \emph{both}), when it exceeds 3 units within a quarter period of
  either end, or when an interior gap is longer than a tenth of the period.
\end{enumerate}
A trace with $\sigma = 0$ is returned whole as baseline (corrected output zero).

\entryfield{Assumptions}
\begin{itemize}
\item \emph{Units.} Not unit-aware (no \code{units} argument), but scale-equivariant by
  construction: multiplying the input by $k$ multiplies the output by $k$, so \emph{sparsity}
  means the same on volts and millivolts.
\item \emph{Sampling rate.} \emph{baseline\_period\_s} is in seconds; the registry allows
  60\,s to 30 days.
\item \emph{Missing values.} Kept as NaN in the output, as above; the count, the trimmed end
  samples and the longest interior gap are reported. Near a long bridged gap the baseline is a
  straight line, not an estimate.
\item \emph{Signal shape.} Events must be the minority and faster than the period; the
  asymmetric penalty assumes they point one way unless \emph{polarity} is \emph{both}.
\item \emph{Span length.} At least 64 samples. A span shorter than the period still runs and
  sets \code{span_shorter_than_period} in the meta.
\end{itemize}

\entryfield{Outputs}
A \texttt{Signal}: the corrected trace, same length and $f_s$. The baseline and the event layer
$x$ are computed but not passed on. \code{meta} carries the solve: \code{fc_cycles_per_sample},
\code{baseline_period_s}, \code{polarity}, \code{r}, \code{lam0}--\code{lam2},
\code{sparsity}, \code{filter_order}, \code{grid_factor}, \code{grid_step_s},
\code{fc_grid_cycles_per_sample}, \code{iterations}, \code{max_iter}, \code{converged},
\code{final_cost}, \code{sigma}, \code{sigma_grid}, \code{sigma_grid_white},
\code{sigma_source}, \code{baseline_min}, \code{baseline_max}, \code{grid_points},
\code{end_pad_grid_points}, \code{end_region_grid_points}, \code{n_nan}, \code{longest_gap_s},
\code{trimmed_end_samples}, \code{event_fraction_3sigma}, \code{span_shorter_than_period},
\code{warnings}, and the adapter's \code{rules}, which the block page reads out.

\entryfield{Cost}
Each iteration is one banded solve, linear in the number of grid points; at most 30 iterations.
Because the solve runs on the grid, a 2.6-million-sample channel becomes about 18\,000 grid
points (module notes; \code{tests/test_beads.py} has a test that a whole channel of that size is
fast and lean). The block registers a calibrated cost model (\code{register_cost_model}) and an
estimate, \code{_estimate}, from the padded grid size; it returns nothing when the inputs would
be refused. Memory is linear in the span. Meant to run locally on whole channels.

\entryfield{Failure modes}
\begin{itemize}
\item \emph{Events as slow as the period} are baseline. The filter's slope is gentle: drift
  at twice the period leaves about 4\% RMS of its amplitude in the output, at 1.5 times about
  12\%, at 3 times about 0.9\% (module notes; bounds checked in
  \code{tests/test_beads.py}).
\item \emph{Event-dominated spans.} If events fill most of the span the baseline settles partway
  into them; levels are then unreliable, shapes and relative depths are kept.
\item \emph{An event near either end} is reflected into the padding and bends the baseline over
  the last half period (measured 0.63 of a drop's depth for a span cut 1040\,s after a drop
  onset); a warning names it.
\item \emph{Drift absorbed as events} at very low \emph{sparsity}: a warning fires above 80\%
  event coverage on a drop or spike run.
\item \emph{Not converged.} After 30 iterations the last iterate is used and
  \code{converged} is false.
\item \emph{Too short a span or period} is refused with a message.
\end{itemize}

\entryfield{Validation}
\code{tests/test_beads.py} holds the evidence. The banded port matches a dense transcription of
\texttt{beads.m} to $10^{-6}$ of the signal scale for orders 1 and 2, and the cost never rises
between iterations. On a synthetic 48\,h, 1\,Hz acceptance trace (6\,h sine drift of amplitude
5, 20 drops of depth 1, noise 0.05) the residual drift away from the drops is below 0.1 RMS,
every drop keeps at least 85\% of its depth, the ends carry no edge artefact above 0.1, and the
overshoot beside the drops is less than half that of a 600\,s rolling mean (measured about
0.06 against 0.34). Other tests cover drift far larger than the noise, AR(0.95) noise, a random
walk, the warnings, NaN handling, scale equivariance and polarity. All of this is on synthetic
traces: no test checks BEADS against hand-marked baselines on a fungal recording.

\entryfield{References}
\citet{ning2014chromatogram}. The reference implementation is Laurent Duval's
\texttt{beads.m} (BSD-3), also adapted by pybaselines; what was verified from the
implementations and what was not is in \code{docs/research/analysis-blocks-primary-sources.md}
(section 6).

\entryfield{Tests}
\code{tests/test_beads.py} (the core, as above); \code{tests/test_adapter_beads.py}
(registration, parameter bounds, that the output is the input minus a baseline, the meta and its
rules, warnings as rules, refusal of a short span, the estimate);
\code{tests/test_adapter_spec.py}, \code{tests/test_block_standard.py} and
\code{tests/test_end_to_end.py} (registry and chain plumbing).
